\scanpage{19}
\begin{center}--- 47 ---\end{center}

\begin{center}\textbf{Algèbre et géométrie}

M. Jean-Pierre \textsc{Serre}, professeur\end{center}

Soit $\mathrm K$ un corps complet pour une valuation discrète, de corps résiduel $k$. On suppose que $\mathrm K$ est de caractéristique zéro et que $k$ est algébriquement clos de caractéristique $p$. Soit $\overline{\mathrm K}$ une clôture algébrique de $\mathrm K$, et soit $\mathrm G$ le groupe de Galois de $\overline{\mathrm K}$ sur $\mathrm K$. Soit $\mathrm V$ un espace vectoriel de dimension finie sur le corps $p$-adique $\mathbf Q_p$, et soit
\[
\varrho:\mathrm G\to\operatorname{Aut}(\mathrm V)
\]
un homomorphisme continu. On dit alors que $\mathrm V$ (ou $\varrho$) est une \textit{représentation $p$-adique} de $\mathrm G$.

Soit $\chi$ l'homomorphisme de $\mathrm G$ dans le groupe des unités $p$-adiques correspondant à l'action naturelle de $\mathrm G$ sur les racines de l'unité d'ordre une puissance de $p$. Soit $\mathrm C$ le complété de $\overline{\mathrm K}$. Si $\mathrm V$ est, comme ci-dessus, une représentation $p$-adique de $\mathrm G$, posons $\mathrm V_{\mathrm C}=\mathrm C\otimes\mathrm V$ (le produit tensoriel étant pris sur $\mathbf Q_p$); pour tout $n\in\mathbf Z$, soit $\mathrm V_{\mathrm C}^n$ le $\mathrm K$-sous-espace de $\mathrm V_{\mathrm C}$ formé des éléments $x$ tels que $\varrho(s)x=\chi(s)^nx$ pour tout $s\in\mathrm G$, et soit $\mathrm V_{\mathrm C}(n)=\mathrm C\otimes_{\mathrm K}\mathrm V_{\mathrm C}^n$. D'après un résultat de \textsc{Tate}, la somme directe des $\mathrm V_{\mathrm C}(n)$ s'identifie à un sous-espace vectoriel de $\mathrm V_{\mathrm C}$. On dit que $\mathrm V$ admet une \textit{décomposition de Hodge $p$-adique} (ou que $\mathrm V$ est un \textit{module de Hodge}) si ce sous-espace vectoriel est égal à $\mathrm V_{\mathrm C}$, autrement dit si l'on a:
\[
\mathrm V_{\mathrm C}=\sum_{n\in\mathbf Z}\mathrm V_{\mathrm C}(n).
\]

Soit $\mathrm V$ un tel module, soit $\mathrm G_{\mathrm V}=\varrho(\mathrm G)$ le sous-groupe correspondant de $\operatorname{Aut}(\mathrm V)$, et soit $\mathrm H_{\mathrm V}$ \textit{l'enveloppe algébrique} de $\mathrm G_{\mathrm V}$, autrement dit le plus petit sous-groupe algébrique (sur $\mathbf Q_p$) du groupe linéaire $\mathbf{GL}_{\mathrm V}$ dont le groupe des points contienne $\mathrm G_{\mathrm V}$.

Le sujet du cours a été l'étude du groupe $\mathrm H_{\mathrm V}$, et en particulier la démonstration des résultats suivants:
\margin{18}
